% ============================================================
%  MTH636 – Game Theory | Lecture 17
%  Subgames, Zermelo's Theorem & Chomp
% ============================================================
\documentclass[12pt, a4paper]{article}

% ============================================================
%  MTH636 -- Game Theory
%  Notes Preamble — Color-Coded, TikZ-Enhanced
% ============================================================

% --- Core packages ---
\usepackage[T1]{fontenc}
\usepackage[utf8]{inputenc}
\usepackage{lmodern}
\usepackage[margin=2.5cm, top=2.8cm, bottom=2.8cm]{geometry}
\usepackage{microtype}

% --- Math ---
\usepackage{amsmath, amssymb, amsthm, mathtools}
\usepackage{bm}        % bold math
\usepackage{cancel}    % strike-through in math

% --- Color & Boxes ---
\usepackage[dvipsnames, svgnames, x11names, table]{xcolor}
\usepackage[most]{tcolorbox}
\tcbuselibrary{breakable, skins, theorems}

% --- TikZ & Diagrams ---
\usepackage{tikz}
\usepackage{tikz-qtree}
\usetikzlibrary{arrows.meta, positioning, shapes.geometric, fit, backgrounds,
                decorations.pathreplacing, calc, matrix, trees}
\usepackage{pgfplots}
\pgfplotsset{compat=1.18}

% --- Tables ---
\usepackage{booktabs}
\usepackage{multirow}
\usepackage{array}
\usepackage{tabularx}
\usepackage{makecell}

% --- Typography & Layout ---
\usepackage{enumitem}
\usepackage{parskip}
\usepackage{fancyhdr}
\usepackage{titlesec}
\usepackage{titletoc}
\usepackage{hyperref}
\usepackage{graphicx}
\usepackage{float}
\usepackage{caption}

% --- Color Palette ---
\definecolor{defblue}{HTML}{1A5276}      % Definition header
\definecolor{defbluebg}{HTML}{D6EAF8}    % Definition background
\definecolor{thmgreen}{HTML}{1E8449}     % Theorem header
\definecolor{thmbg}{HTML}{D5F5E3}        % Theorem background
\definecolor{exorange}{HTML}{B7770D}     % Example header
\definecolor{exbg}{HTML}{FEF9E7}         % Example background
\definecolor{warnred}{HTML}{922B21}      % Warning/Key Insight header
\definecolor{warnbg}{HTML}{FADBD8}       % Warning background
\definecolor{noticegray}{HTML}{2C3E50}   % Note header
\definecolor{noticebg}{HTML}{EBF5FB}     % Note background
\definecolor{taketawayyellow}{HTML}{7D6608} % Key Takeaway header
\definecolor{takeawaybg}{HTML}{FDFEFE}   % Takeaway background
\definecolor{accentpurple}{HTML}{6C3483} % Accent / proof
\definecolor{lightgray}{HTML}{F2F3F4}    % shading

% ============================================================
%  TCOLORBOX ENVIRONMENTS
% ============================================================

% -- Definition (Blue) --
\newtcolorbox{defbox}[2][]{
  enhanced, breakable,
  colback=defbluebg, colframe=defblue, coltitle=white,
  fonttitle=\bfseries\sffamily,
  title={Definition: #2},
  attach boxed title to top left={yshift=-2mm, xshift=4mm},
  boxed title style={colback=defblue, rounded corners},
  arc=4pt, boxrule=1.2pt,
  left=6pt, right=6pt, top=8pt, bottom=6pt,
  #1
}

% -- Theorem / Proposition (Green) --
\newtcolorbox{thmbox}[2][]{
  enhanced, breakable,
  colback=thmbg, colframe=thmgreen, coltitle=white,
  fonttitle=\bfseries\sffamily,
  title={Theorem: #2},
  attach boxed title to top left={yshift=-2mm, xshift=4mm},
  boxed title style={colback=thmgreen, rounded corners},
  arc=4pt, boxrule=1.2pt,
  left=6pt, right=6pt, top=8pt, bottom=6pt,
  #1
}

% -- Proposition (Green variant) --
\newtcolorbox{propbox}[2][]{
  enhanced, breakable,
  colback=thmbg, colframe=thmgreen!70!black, coltitle=white,
  fonttitle=\bfseries\sffamily,
  title={Proposition: #2},
  attach boxed title to top left={yshift=-2mm, xshift=4mm},
  boxed title style={colback=thmgreen!70!black, rounded corners},
  arc=4pt, boxrule=1.2pt,
  left=6pt, right=6pt, top=8pt, bottom=6pt,
  #1
}

% -- Example (Orange/Yellow) --
\newtcolorbox{exbox}[2][]{
  enhanced, breakable,
  colback=exbg, colframe=exorange, coltitle=white,
  fonttitle=\bfseries\sffamily,
  title={Example: #2},
  attach boxed title to top left={yshift=-2mm, xshift=4mm},
  boxed title style={colback=exorange, rounded corners},
  arc=4pt, boxrule=1.2pt,
  left=6pt, right=6pt, top=8pt, bottom=6pt,
  #1
}

% -- Key Insight / Warning (Red) --
\newtcolorbox{warnbox}[2][]{
  enhanced, breakable,
  colback=warnbg, colframe=warnred, coltitle=white,
  fonttitle=\bfseries\sffamily,
  title={Key Insight: #2},
  attach boxed title to top left={yshift=-2mm, xshift=4mm},
  boxed title style={colback=warnred, rounded corners},
  arc=4pt, boxrule=1.2pt,
  left=6pt, right=6pt, top=8pt, bottom=6pt,
  #1
}

% -- Note (Gray-Blue) --
\newtcolorbox{notebox}[1][]{
  enhanced, breakable,
  colback=noticebg, colframe=noticegray!60, coltitle=white,
  fonttitle=\bfseries\sffamily,
  title={Note},
  attach boxed title to top left={yshift=-2mm, xshift=4mm},
  boxed title style={colback=noticegray!60, rounded corners},
  arc=4pt, boxrule=0.8pt,
  left=6pt, right=6pt, top=8pt, bottom=6pt,
  #1
}

% -- Key Takeaway --
\newtcolorbox{takebox}[1][]{
  enhanced, breakable,
  colback=takeawaybg, colframe=taketawayyellow, coltitle=white,
  fonttitle=\bfseries\sffamily,
  title={Key Takeaway},
  attach boxed title to top left={yshift=-2mm, xshift=4mm},
  boxed title style={colback=taketawayyellow, rounded corners},
  arc=4pt, boxrule=1.2pt,
  left=6pt, right=6pt, top=8pt, bottom=6pt,
  #1
}

% -- Intuition --
\newtcolorbox{intubox}[1][]{
  enhanced, breakable,
  colback=accentpurple!8, colframe=accentpurple!60, coltitle=white,
  fonttitle=\bfseries\sffamily,
  title={Intuition},
  attach boxed title to top left={yshift=-2mm, xshift=4mm},
  boxed title style={colback=accentpurple!80, rounded corners},
  arc=4pt, boxrule=0.8pt,
  left=6pt, right=6pt, top=8pt, bottom=6pt,
  #1
}

% ============================================================
%  PAYOFF MATRIX HELPER
%  Usage: inside tabular, use \payoff{a}{b} for (a,b)
% ============================================================
\newcommand{\payoff}[2]{$(#1,\;#2)$}
\newcommand{\nash}[2]{\underline{$#1$},\underline{$#2$}}  % underline NE payoffs
\newcommand{\Nset}{\mathbb{N}}
\newcommand{\Rset}{\mathbb{R}}
\newcommand{\Sset}{\mathcal{S}}
\newcommand{\Ui}{u_i}

% ============================================================
%  SECTION FORMATTING
% ============================================================
\titleformat{\section}
  {\large\bfseries\sffamily\color{defblue}}
  {\thesection}{1em}{}[\titlerule]

\titleformat{\subsection}
  {\normalsize\bfseries\sffamily\color{thmgreen!80!black}}
  {\thesubsection}{1em}{}

\titleformat{\subsubsection}
  {\normalsize\itshape\sffamily\color{exorange!80!black}}
  {\thesubsubsection}{1em}{}

% ============================================================
%  HEADER / FOOTER
% ============================================================
\pagestyle{fancy}
\fancyhf{}
\fancyhead[L]{\sffamily\small\color{defblue} MTH636 -- Game Theory}
\fancyhead[R]{\sffamily\small\color{defblue} \leftmark}
\fancyfoot[C]{\sffamily\small\thepage}
\renewcommand{\headrulewidth}{0.4pt}

% ============================================================
%  HYPERREF SETUP
% ============================================================
\hypersetup{
  colorlinks=true,
  linkcolor=defblue,
  urlcolor=thmgreen,
  citecolor=accentpurple,
  bookmarksnumbered=true,
  pdftitle={MTH636 Game Theory Notes},
  pdfauthor={MTH636 Student Notes}
}

% ============================================================
%  MISC HELPERS
% ============================================================
\newcommand{\BR}{\text{BR}}        % Best Response
\newcommand{\NE}{\text{NE}}        % Nash Equilibrium
\newcommand{\SPNE}{\text{SPNE}}    % Subgame Perfect NE
\newcommand{\Dom}{\text{dom}}      % Dominance
\newcommand{\mixs}[1]{\sigma_{#1}} % Mixed strategy
\newcommand{\pures}[1]{s_{#1}}     % Pure strategy


% ============================================================
\begin{document}

% Title
\begin{center}
  {\LARGE\bfseries\sffamily\color{defblue} MTH636 -- Game Theory}\\[6pt]
  {\Large\sffamily Lecture 17: Subgames \& Zermelo's Theorem}\\[4pt]
\end{center}
\vspace{0.5cm}
\hrule height 1.5pt
\vspace{0.3cm}

% ============================================================
\section{Strategy Profiles and Paths}

In an extensive form game $\Gamma = \langle N, V, E, n_0, (V_i)_{i \in N}, (u_i)_{i \in N} \rangle$, a strategy for player $i$ is a map $s_i: V_i \to \cup A(n)$.

A \textbf{strategy profile} $s = (s_1, \dots, s_n)$ is a tuple containing a strategy for each player. 
Every strategy profile uniquely determines a specific path in the game tree, starting from the root $n_0$ and ending at a leaf node $n_k$:
\[ (n_0, n_1, \dots, n_k) \]
The utility of the strategy profile is simply defined as the utility at the leaf node reached:
\[ u(s) := u(n_k) \]

% ============================================================
\section{Subgames}

\begin{defbox}{Subgame}
  Let $\Gamma$ be an extensive form game. Consider a node $n \in V$. The \textbf{subgame starting at $n$}, denoted by $\Gamma(n)$, is an extensive form game itself:
  \[ \Gamma(n) = \langle N, V(n), E(n), n, (V_i(n))_{i \in N}, (u_i)_{i \in N} \rangle \]
  where:
  \begin{itemize}
      \item $N$: The set of players remains the same.
      \item $V(n)$: The set of all nodes that can be reached from $n$ (descendants of $n$, including $n$).
      \item $E(n)$: All the edges that connect nodes within $V(n)$.
      \item $n$: The new root of the subgame.
      \item $V_i(n) = V_i \cap V(n)$: The nodes controlled by player $i$ within the subgame.
      \item $u_i$: The restriction of the utility function $u_i$ to the leaves in the new game tree $(V(n), E(n), n)$.
  \end{itemize}
\end{defbox}

\textbf{Remarks:}
\begin{itemize}
    \item The original game is a subgame of itself: $\Gamma = \Gamma(n_0)$.
    \item Every leaf node trivially defines a subgame where no player makes a move, and payoffs are immediately realized.
\end{itemize}

% ============================================================
\section{Win/Loss/Draw Games \& Zermelo's Theorem}

Consider a particular type of game:
\begin{itemize}
    \item Two players.
    \item Perfect information.
    \item Finite duration.
    \item Only three possible outcomes:
    \begin{enumerate}
        \item Player 1 wins (payoffs $1, -1$)
        \item Player 2 wins (payoffs $-1, 1$)
        \item The game ends in a draw (payoffs $0, 0$)
    \end{enumerate}
\end{itemize}
Examples include Tic-Tac-Toe, Checkers, and Chess (ignoring infinite loops).

\begin{defbox}{Winning and Drawing Strategies}
  \begin{itemize}
      \item A \textbf{winning strategy} for Player 1 is a strategy $s_1$ such that $u(s_1, s_2) = (1, -1)$ for \textit{all} possible opponent strategies $s_2 \in S_2$.
      \item A \textbf{strategy guaranteeing at least a draw} for Player 1 is a strategy $s_1$ such that $u(s_1, s_2) \in \{(1,-1), (0,0)\}$ for all $s_2 \in S_2$.
  \end{itemize}
  (Analogous definitions apply for Player 2).
\end{defbox}

\begin{thmbox}{Zermelo's Theorem (1913) / von Neumann (1928)}
  In every two-player, finite game of perfect information in which the set of utilities is $\{(1,-1), (-1,1), (0,0)\}$, \textbf{one and only one} of the following holds:
  \begin{enumerate}
      \item Player 1 has a winning strategy.
      \item Player 2 has a winning strategy.
      \item Both players have strategies guaranteeing at least a draw.
  \end{enumerate}
\end{thmbox}

% ============================================================
\section{Introduction to Chomp (David Gale, 1974)}

To illustrate these concepts, we introduce the game \textbf{Chomp}.
\begin{itemize}
    \item Played on an $m \times n$ grid of squares (like a chocolate bar).
    \item The bottom-left square is "poisoned."
    \item Players take turns choosing a valid square $(i,j)$. 
    \item When a square is chosen, that square and all squares above and to the right of it are "eaten" and removed from the board.
    \item The player forced to eat the poisoned bottom-left square loses.
\end{itemize}
Because Chomp is finite, has perfect information, and has no draws, Zermelo's theorem guarantees that either Player 1 or Player 2 must have a winning strategy. We will prove who wins in the next lecture.

\end{document}
